\begin{obeassessment}
\textbf{Kuis Bab 16}
1. Sebutkan empat layanan kriptografi dan jelaskan masing-masing dalam satu
   kalimat.
2. Jelaskan mengapa nir-penyangkalan berbeda dari otentikasi pengirim.
3. Jelaskan peran kunci publik dan kunci privat pada tanda tangan digital.
4. Sebutkan algoritma tanda tangan digital yang disebut modul dan standar yang
   mengatur tanda tangan RSA.
\end{obeassessment}
\begin{competencychecklist}
\begin{tabularx}{\linewidth}{|X|c|}
\hline
Indikator Kompetensi & Tercapai \\
\hline
Saya dapat menyebutkan empat layanan kriptografi & [ ] \\
\hline
Saya dapat menjelaskan makna nir-penyangkalan & [ ] \\
\hline
Saya dapat menjelaskan peran kunci publik dan kunci privat pada tanda tangan digital & [ ] \\
\hline
Saya dapat mengidentifikasi algoritma tanda tangan digital yang disebut modul & [ ] \\
\hline
\end{tabularx}
\end{competencychecklist}
